\begin{abstract}
Reversible data hiding returns the cover exactly, but its hidden bits are lost when the stego image is changed; deep
steganography networks are robust, but floating-point rounding stops them from returning the cover. We
propose INN-RC, a lossless steganography scheme built on a message-conditioned integer coupling network (IC-INN). The
cover is split into two pixel halves that are updated in turn by rounded offsets predicted from the other half and the
message, so the receiver can undo the network exactly once it knows the message and no latent variable has to be
stored. A small decoder reads the message back. The bits it gets wrong and the pixels clipped at 0 or 255 are coded with
a context-adaptive arithmetic coder, whose contexts are the decoder confidence and the coupling offset re-estimated by the
receiver, and are hidden with multi-layer rhombus prediction-error expansion. The 62.9k-parameter model trains on a CPU and the
pipeline runs in MATLAB. On \nImgs{} test images from BSD100, Kodak and real DICOM CT/MR/CR scans,
every successfully embedded image was recovered bit for bit with zero bit error rate. At 0.0156 bpp (1024 bits) the
stego PSNR is \psnrQfour{} dB (SSIM \ssimQfour{}) and the side information averages \auxQfour{} bits, about half the
message length, while a float invertible network with the same compensation would need \auxBASE{} bits. Unlike classical
prediction-error expansion, whose bits are lost under any distortion, the message survives Gaussian noise ($\sigma=2$)
and 1\% salt-and-pepper noise with a BER of \robgausstwoR\% and \robspsponeR\%.
\end{abstract}
